\documentclass{article}

% if you need to pass options to natbib, use, e.g.:
%     \PassOptionsToPackage{numbers, compress}{natbib}
% before loading neurips_2026

% The authors should use one of these tracks.
% Before accepting by the NeurIPS conference, select one of the options below.
% 0. "default" for submission
\usepackage{neurips_2026}
% the "default" option is equal to the "main" option, which is used for the Main Track with double-blind reviewing.
% 1. "main" option is used for the Main Track
%  \usepackage[main]{neurips_2026}
% 2. "position" option is used for the Position Paper Track
%  \usepackage[position]{neurips_2026}
% 3. "eandd" option is used for the Evaluations & Datasets Track
 % \usepackage[eandd]{neurips_2026}
% 4. "creativeai" option is used for the Creative AI Track
%  \usepackage[creativeai]{neurips_2026}
% 5. "sglblindworkshop" option is used for the Workshop with single-blind reviewing
 % \usepackage[sglblindworkshop]{neurips_2026}
% 6. "dblblindworkshop" option is used for the Workshop with double-blind reviewing
%  \usepackage[dblblindworkshop]{neurips_2026}
% 7. "education" option is used for the Education Track
%  \usepackage[education]{neurips_2026}


% After being accepted, the authors should add "final" behind the track to compile a camera-ready version.
% 1. Main Track
 % \usepackage[main, final]{neurips_2026}
% 2. Position Paper Track
%  \usepackage[position, final]{neurips_2026}
% 3. Evaluations & Datasets Track
 % \usepackage[eandd, final]{neurips_2026}
% 4. Creative AI Track
%  \usepackage[creativeai, final]{neurips_2026}
% 5. Workshop with single-blind reviewing
%  \usepackage[sglblindworkshop, final]{neurips_2026}
% 6. Workshop with double-blind reviewing
%  \usepackage[dblblindworkshop, final]{neurips_2026}
% Note. For the workshop paper template, both \title{} and \workshoptitle{} are required, with the former indicating the paper title shown in the title and the latter indicating the workshop title displayed in the footnote.
% For workshops (5., 6.), the authors should add the name of the workshop, "\workshoptitle" command is used to set the workshop title.
% \workshoptitle{WORKSHOP TITLE}
% 7. Education Track
% \usepackage[education, final]{neurips_2026}

% "preprint" option is used for arXiv or other preprint submissions
 % \usepackage[preprint]{neurips_2026}

% to avoid loading the natbib package, add option nonatbib:
%    \usepackage[nonatbib]{neurips_2026}

\usepackage[utf8]{inputenc} % allow utf-8 input
\usepackage[T1]{fontenc}    % use 8-bit T1 fonts
\usepackage{hyperref}       % hyperlinks
\usepackage{url}            % simple URL typesetting
\usepackage{booktabs}       % professional-quality tables
\usepackage{amsfonts}       % blackboard math symbols
\usepackage{nicefrac}       % compact symbols for 1/2, etc.
\usepackage{microtype}      % microtypography
\usepackage{xcolor}         % colors
\usepackage{amsmath,amssymb,amsthm}
\usepackage{graphicx}
\graphicspath{{../figures/}}
\newtheorem{proposition}{Proposition}
\newtheorem{definition}{Definition}

% Note. For the workshop paper template, both \title{} and \workshoptitle{} are required, with the former indicating the paper title shown in the title and the latter indicating the workshop title displayed in the footnote. 
\title{How Much Targeted Training Does It Take to Top a Multi-Task Leaderboard? Minimum Investment Under Different Aggregation Rules}

\author{Anonymous Authors}

\begin{document}
\maketitle

\begin{abstract}
A leaderboard turns per-task scores into one rank, so a developer can try to lift a chosen model by training on data related to individual sub-datasets. How much such targeted training is needed to put the model first, and how does that depend on the aggregation rule? We define the investment as the minimum raise of the target's task scores, within a per-task cap, that makes it strictly first, and measure it in points, tasks touched and largest single-task jump. We solve it exactly for score rules (mean, median, minimum, power means); for rank rules (Borda, Condorcet winner, plurality; Copeland on HELM) an integer program under strict comparisons decides feasibility and tasks touched exactly and returns the points as a witness cost within $n\varepsilon$ of the infimum. We prove closed forms for the mean, minimum and odd median, and apply it to every model of the Open LLM Leaderboard (OLLB; 4495 targets, 6 tasks) and of HELM MMLU (78 targets, 57 subjects). On OLLB the rank rules are solved only for the 49 best non-leading targets (mean-ranks 2--50), so comparisons involving them use those targets alone. With caps taken from observed fine-tuning gains, first place is reachable for 12\% of OLLB targets under the mean (bootstrap interval 11--17\%; every target up to rank 394 is feasible; the largest feasible rank is 901) and for 19\% of HELM targets (15--42\%; up to rank 16); with a 100-point cap every target can reach it. The rules do not yield one ordering. In points, on targets feasible with positive cost under both rules, the median needs a median of 0.39 of the mean's total on OLLB (0.07 on HELM) and Borda 0.27 (0.51), but not at the very top: the OLLB targets of mean-rank 2--4 need more points under the median than under the mean. In reach, the OLLB median makes fewer targets feasible than the mean (10\% against 12\%, all-feasible prefix 90 against 394) although its largest feasible rank is larger (1207 against 901; the two shares' intervals overlap), while on HELM its point estimate is higher (33\% against 19\%, intervals overlapping). Conclusions are conditional on the cap, which on HELM is transferred from OLLB, and hold for fixed rivals and fixed score tables.
\end{abstract}

\section{Introduction}
\label{sec:intro}
Multi-task benchmarks rank models by aggregating per-task scores into one number, and the resulting rank influences which models are adopted, published and funded. A developer therefore has an incentive to improve a chosen target model specifically on the benchmark's sub-datasets, for instance by training on data related to them. How much such targeted improvement is needed to reach first place, and how that amount depends on the aggregation rule, is the question we study on public leaderboard data.

Our approach is deliberately simple. We take the score table of a public leaderboard as fixed, let the target raise its own scores task by task up to a cap, and compute the \emph{minimum} raise that makes it strictly first under a given aggregation rule. The cap is the central modelling assumption: it stands for how far training on related data can move one score. We therefore report results under data-grounded caps and an oracle cap, and report infeasibility (no allowed raise makes the target first) as an outcome rather than a missing value.

\paragraph{Contributions.}
(i) A formulation of the minimum targeted investment with three cost measures ($L_1$, $K$, $L_\infty$) and an explicit family of caps (Section~\ref{sec:setup}). (ii) Exact solvers for score rules (mean, median, minimum, geometric, harmonic and squared power means) and integer programs for rank rules (Borda, Condorcet winner, plurality, Copeland) that decide feasibility and $K$ exactly and return an $L_1$ witness cost within $n\varepsilon$ of the infimum, all tested against brute force, with closed forms and proofs for the mean, minimum and odd-$n$ median (Section~\ref{sec:solvers}, Appendix~\ref{app:proofs}). (iii) A study of every model on two public leaderboards, with bootstrap and evaluation-noise ranges, paired same-target rule comparisons and a feasibility frontier (Sections~\ref{sec:exp}--\ref{sec:results}); the rank rules on OLLB cover only the 49 best non-leading targets, as stated wherever they appear.

\paragraph{Findings.}
Under caps grounded in observed fine-tuning gains, first place is out of reach for most models under every rule studied: the feasible fraction under the mean is 12\% on OLLB and 19\% on HELM. The rules differ in two ways that do not combine into one ordering. First, \emph{cost on the same targets}: on targets feasible with positive cost under both rules, order-statistic and rank rules need a fraction of the mean's total points (median-of-ratios on OLLB: median 0.39, minimum 0.11, Borda 0.27, Condorcet 0.52; the Borda, Condorcet and plurality ratios use the 36--48 best targets only), with exceptions at the very top, where the OLLB median costs more than the mean for the targets of mean-rank 2--4. Second, \emph{reach}: the OLLB median is feasible for fewer targets than the mean (10\% against 12\%) and has a shorter all-feasible prefix (90 against 394) but a larger largest feasible rank (1207 against 901); on HELM the median's point estimate is 33\% against 19\% for the mean, with overlapping bootstrap intervals. Of three hypotheses fixed beforehand, one holds, one holds in part and one is not supported for the median rule (Section~\ref{sec:results}). All conclusions are conditional on the cap and on the other models being fixed; Section~\ref{sec:disc} lists where they fail.

\section{Related work}
\label{sec:related}
\paragraph{Benchmarks and leaderboards.}
MMLU \citep{hendrycks2021mmlu} aggregates 57 subjects, and HELM \citep{liang2023helm} releases per-scenario scores and raw model outputs for many models; BIG-Bench Hard \citep{suzgun2023bbh} selects hard tasks from BIG-Bench. These provide the multi-task score matrices that we analyse.

\paragraph{Fragility of rankings.}
\citet{dehghani2021lottery} show that the relative order of methods can change with the choice of benchmark tasks (the ``benchmark lottery''). \citet{singh2025illusion} document how private testing and unequal data access distort the Chatbot Arena leaderboard. Both concern what is measured and disclosed; we instead hold the leaderboard fixed and ask what a single developer must change for one model to win.

\paragraph{Aggregation and social choice.}
\citet{lanctot2025vase} treat each task as a voter and use voting theory to build cross-task evaluation; the Handbook of Computational Social Choice \citep{brandt2016handbook} covers the rules (Borda, Condorcet-consistent rules, plurality) we reuse. Our rules are the same families, but our question is the cost of changing the winner, not the choice of rule.

\paragraph{Training on the test task and contamination.}
\citet{zhou2023cheater} show that benchmark leakage into training can raise scores sharply, and \citet{yang2023rephrased} show that rephrased test samples evade common decontamination. \citet{chen2026incentives} analyse, in a game-theoretic model, the incentives that benchmarks create for strategic post-training. Our analysis is empirical and combinatorial: given a fixed score table, we compute the exact minimum targeted improvement per rule, under explicit assumptions about how much any one task can be improved.

\paragraph{Positioning.}
Closest in question is \citet{chen2026incentives} (we read the full text of arXiv v1). Their developers choose a scalar amount of benchmark-specific effort that a post-effort score maps to one benchmark score per model; a multi-dimensional capability with a linear benchmark weight vector is mentioned as an extension, but the rule that aggregates tasks is not a variable of their analysis. They prove that current protocols induce games without a Nash equilibrium and that tune-before-test induces a unique equilibrium ranking by latent quality, and in a case study of one benchmark they measure the training steps needed to change a ranking under tune-before-test. We take the opposite stance: no game and no equilibrium, a fixed table, and the exact cost for one developer under each of several aggregation rules, in score points under per-task caps rather than training steps. \citet{colombo2022best} propose a social-choice-based aggregation for NLP benchmarks and show it changes conclusions relative to the mean; they also corrupt the scores of some tasks for all systems at random and find rank-based aggregation more stable than the mean in ranking error. We do not propose a rule, and our perturbation is different: one target raises its own scores, and we compute the minimum cost of becoming first under each rule. The computational-social-choice literature studies the analogous \emph{bribery} problem, the cost of changing votes so that a chosen candidate wins under a voting rule \citep{faliszewski2009bribery,brandt2016handbook}. Our setting differs in that the manipulator changes its own scores on tasks (the ``voters'') and pays in score points under per-task caps, which gives closed forms for score rules and integer programs for rank rules.

\section{Problem setup}
\label{sec:setup}
\subsection{Leaderboard, target and investment}
A leaderboard is a score matrix $X\in[0,100]^{M\times n}$ of $M$ models on $n$ tasks. Fix a target model $i$ with score vector $x=X_{i\cdot}$; all other rows stay fixed. In targeted training the developer can raise the target's score on task $t$ to any $y_t\in[x_t,c_t]$, where $c_t$ is a \emph{cap} (Section~\ref{sec:caps}). An aggregation rule $f$ maps the score vectors to a ranking; the target is \emph{first} if it is strictly ahead under $f$. For a score rule we require a margin $\varepsilon=10^{-3}$ points, $f(y)\ge\tau+\varepsilon$ with $\tau$ the best rival's value. A rank rule uses plain strict comparisons and no margin: a pairwise comparison of two scores is a win (strictly higher), a loss, or a tie worth one half, and under Borda, plurality and Copeland the target is first if its total is strictly the largest, under the Condorcet requirement if it beats every rival on more than half of the tasks; there $\varepsilon$ only sets the granularity of the witness we return. For score rules with $\varepsilon>0$ the minima below are attained, and the infimum over strict wins is the limit $\varepsilon\to0$; for rank rules the infimum over strict wins need not be attained (it can be attained at a tie), so the solver returns an explicit witness (Section~\ref{sec:solvers}). Evaluation noise of the scores is far larger than $\varepsilon$ (Section~\ref{sec:results}), so the score margin defines an exact optimisation problem, not a win that would survive resampling; for score rules, margins of 1 and 2 points are reported as variants. We report the investment in three ways:
\[
L_1=\sum_t (y_t-x_t),\qquad K=\#\{t: y_t>x_t\},\qquad L_\infty=\max_t (y_t-x_t),
\]
the total points, the number of tasks touched, and the largest single-task jump. The minimum of each is taken separately, so the three need not come from one raise. $L_\infty$ is the smallest $d$ for which a raise of every task by at most $d$ (a ``general improvement'') suffices. For comparison across rules and datasets we also report $K/n$ and $L_1/n$.

\subsection{Caps: how much can one task be improved?}
\label{sec:caps}
Without a bound on $y_t$ the question is trivial (put every task at 100). The cap $c_t$ is therefore the central modelling assumption, and we never collapse it to a single choice. The primary caps are data-grounded: on OLLB \texttt{emp99}, $c_t=\min(100,x_t+g_t^{99})$ where $g_t^{99}$ is the 99th percentile of the observed gains (fine-tune minus its declared base model, 540 pairs) on task $t$; on HELM, which has no such pairs, \texttt{emp99T}, the OLLB MMLU-PRO value $g^{99}_{\text{MMLU-PRO}}$ applied to every subject (a transfer assumption). Both inherit a selection bias, since only published fine-tunes appear, they measure gains attained by the community and not those a determined developer could attain, and the 99th percentile of 540 pairs rests on about five tail points per task. The oracle cap \texttt{ceil}, $c_t=100$, gives a lower bound on any required investment. Further caps (95th percentiles, the best rival's score, and $+5/+10/+20$ points on every task) are reported in the appendix. When no $y\le c$ makes the target first, the case is \emph{infeasible}; infeasibility is reported as a rate and never coded as a large investment, and all investment statistics are conditional on feasibility.

\subsection{Aggregation rules}
\label{sec:rules}
Score rules apply a function $f$ to the score vector of each model: the \emph{mean} (the rule of both leaderboards we use), the \emph{median}, and the \emph{minimum} (worst-task score). Rank rules use only the order of models on each task: \emph{Borda} (sum over tasks of the number of models beaten, ties counting one half; equivalent to mean rank), the \emph{Condorcet-winner} requirement (the target beats every rival on a strict majority of tasks, ties counting one half), \emph{Copeland} (a model scores 1 for each rival it beats by that pairwise count, one half for a tie; first means the strictly largest score) and \emph{plurality} (each task gives one point, split equally among the models tied for its highest score; first means the strictly largest total). Being a Condorcet winner is sufficient for being first under any Condorcet-consistent rule such as Copeland or Schulze, so its investment upper-bounds theirs. Geometric, harmonic and power means and the trimmed mean are in the appendix.

\subsection{Exact minimum-investment solvers}
\label{sec:solvers}
For a score rule let $\tau=\max_{j\neq i} f(X_{j\cdot})$ be the level of the best rival; the target is first iff $f(y)\ge \tau+\varepsilon$. For the mean, the minimum and the median the minimum investment has a closed form (Propositions~\ref{prop:mean}--\ref{prop:median}, proved in Appendix~\ref{app:proofs}); for power means and the trimmed mean we solve a convex program or enumerate exactly. For rank rules the outcome depends only on the order of scores, so we use plain strict comparisons and no margin. The integer program (HiGHS, rival constraints added lazily) lets each raised task end at a rival's score (a tie) or at most $\varepsilon$ above it, and returns that explicit witness. By construction and by tests against brute force on fractional, tied and cap-boundary instances, feasibility and $K$ are exact for continuous raises, and the true infimum of $L_1$ lies within $n\varepsilon$ of the witness cost (Appendix~\ref{app:ordinal}). A target already ahead of every rival under a score rule is counted first at zero cost, whatever the size of its lead. All solvers are tested against brute force on small instances. For each target we report $L_1$ and $K$, each minimised separately; $L_\infty$ is computed but is not part of the headline results.

Write $b=\tau+\varepsilon$ for the level a score rule must reach and $h_t=c_t-x_t$ for the headroom of task $t$. They concern targets not already ahead ($f(x)\le\tau$); a target with $f(x)>\tau$ costs zero by the convention above.
\begin{proposition}[Mean]\label{prop:mean}
Let $f$ be the mean and $R=n\,b-\sum_t x_t$ the required total. If $R\le0$ the target is already first and every cost is zero. Otherwise it can become first iff $\sum_t h_t\ge R$; the minimum $L_1$ is $R$ whatever its spread; the minimum $K$ is the least $k$ for which the $k$ largest $h_t$ sum to at least $R$; the minimum $L_\infty$ is the least $d$ with $\sum_t\min(d,h_t)\ge R$, which equals $R/n$ when no cap binds. The capped uniform raise $y_t=\min(x_t+d,c_t)$ with this $d$ attains the minimum $L_1$ and the minimum $L_\infty$.
\end{proposition}
\begin{proposition}[Minimum]\label{prop:min}
Let $f=\min_t y_t$ and $D=\{t:x_t<b\}$. The target can become first iff $c_t\ge b$ for all $t\in D$. Then $L_1=\sum_{t\in D}(b-x_t)$ and $K=|D|$: every deficient task must be raised, and $L_\infty=\max_{t\in D}(b-x_t)$ (zero if $D$ is empty).
\end{proposition}
\begin{proposition}[Median, $n$ odd]\label{prop:median}
Let $m=(n+1)/2$, $a=\#\{t:x_t\ge b\}$ and, for $x_t<b$, $d_t=b-x_t$ if $c_t\ge b$ and $d_t=\infty$ otherwise. The target is first iff at least $m$ tasks end at or above $b$. If $a\ge m$ no raise is needed; otherwise the minimum $L_1$ is the sum of the $m-a$ smallest finite $d_t$ and the minimum $K$ is $m-a$, and the target can be made first iff $m-a$ finite $d_t$ exist.
\end{proposition}
These are finite-margin statements: the minima are attained at $b=\tau+\varepsilon$. For the strict-win infimum let $\varepsilon\to0$, which replaces $b$ by $\tau$ and every ``$\ge b$'' by ``$>\tau$'' in the feasibility conditions; the infimum of $L_1$ is then $R_0=n\tau-\sum_tx_t$ for the mean and $\sum_{t\in D_0}(\tau-x_t)$ for the minimum, but for a target that is not already ahead it is not attained (a strict win needs a strictly positive margin). Near ties between a rival's level and a task score the finite-margin $K$ can differ from the infimum's $K$ by the tasks within $\varepsilon$ of $\tau$; the code uses $\varepsilon=10^{-3}$ points throughout. For even $n$ the median is the average of the two middle order statistics and we only have the computed solution. We verify Propositions~\ref{prop:mean}--\ref{prop:median} numerically against the solvers (Section~\ref{sec:exp}); no theorem is claimed for the rank rules, which follow the strict model above (Appendix~\ref{app:ordinal}).

\section{Experimental setup}
\label{sec:exp}
\paragraph{Data.}
We use two public leaderboards, each at a fixed, recorded version (checksums in the supplementary reproduction notes). \emph{HELM MMLU}, release v1.13.0 (summary dated 2025-01-10): the exact-match score of 79 models on each of the 57 MMLU subjects ($79\times57$, no missing cell). \emph{Open LLM Leaderboard} (OLLB) \texttt{contents}, dataset commit \texttt{9c09a7c} (2025-03-20): 4576 rows, of which 4497 have distinct model names after we keep the best-scoring variant of each duplicated name, and 4496 after removing one flagged model; each model has six normalised task scores (IFEval, BBH, MATH Lvl 5, GPQA, MUSR, MMLU-PRO), so $n=6$. Merged models are kept in the main analysis and dropped in a robustness slice. The mean we use equals the sites' own aggregate (OLLB \texttt{Average}, HELM ``MMLU All Subjects'') to within $6\times10^{-14}$ points on every model. Neither source declares a licence for the score data; we redistribute only the derived score tables.

\paragraph{Targets, caps and what is reported.}
Every model other than the current leader is a target in turn, with all others fixed (4495 targets on OLLB, 78 on HELM). On OLLB the primary cap is \texttt{emp99} (per-task 99th percentile of the gain of a fine-tune over its declared base model, estimated from 540 such pairs in the same leaderboard; this is about five observations in each task's upper tail; its sampling error is quantified below by a bootstrap over the pairs). HELM has no base/fine-tune pairs, so its primary cap, \texttt{emp99T}, applies the OLLB MMLU-PRO 99th-percentile gain to all 57 subjects; this is a transfer assumption, not a measurement on HELM. The oracle cap \texttt{ceil} and the further caps of Appendix~\ref{app:more} are reported beside it. All investment statistics are medians over feasible targets, shown together with the feasible fraction. Rules are compared on the same targets only (paired comparisons, with the number of targets in both feasible sets). OLLB rank rules are integer programs that at $M=4496$ models we run only for the 50 best targets under the mean (49 non-leading targets; mean-ranks 2--50, a 120\,s limit per instance, time-outs reported as censored, never as infeasible), plus a variant on the 500 best models (Appendix~\ref{app:more}); Copeland is solved on HELM and on the 500-best OLLB slice.

\section{Results}
\label{sec:results}
All statements below are about targets that can be made first within the cap; the feasible fraction is shown beside every cost. A rule is called cheaper than another only on targets feasible under both, and any OLLB statement that involves Borda, Condorcet or plurality refers to the 49 best non-leading targets.

\paragraph{The closed forms are exact (Q1).}
Propositions~\ref{prop:mean}--\ref{prop:median} agree with the exact solvers on every target of both datasets and every cap: the largest deviation in $L_1$ is below $2\times10^{-12}$ points, the deviation in $K$ is zero, and no feasibility verdict differs (Appendix Table~\ref{tab:T2}). The OLLB median is solver-only because $n=6$ is even.

\begin{figure}[t]
\centering
\includegraphics[width=\linewidth]{F_ollb_primary.pdf}\\[2pt]
\includegraphics[width=\linewidth]{F_helm_primary.pdf}
\caption{Minimum targeted investment to become first versus the target's current rank under the mean, primary caps. Left: fraction of targets that can reach first at all. Middle and right: tasks touched ($K/n$) and total points per task ($L_1/n$), \emph{among feasible targets}. Top: OLLB, $n=6$, cap \texttt{emp99} (4495 targets; Borda and Condorcet only for the 49 best non-leading targets, ranks 2--50). Bottom: HELM MMLU, $n=57$, cap \texttt{emp99T} (78 targets; the cap is transferred from OLLB). Points are bin medians (bins double from rank 2; bins with fewer than 3 targets omitted), bands are 95\% bootstrap intervals over targets. The feasible sets differ by rule, so vertical gaps between rules in the middle and right panels compare different targets; same-target comparisons are in Table~\ref{tab:T3}. Generated by \texttt{figures/src/fig\_rank\_rules.py}.}
\label{fig:main}
\end{figure}

\paragraph{Feasibility depends on rank, rule and cap (Q3).}
Under the data-grounded caps, targeted training cannot make most models first (Figure~\ref{fig:main}, left). Appendix Table~\ref{tab:T4} gives three summaries that do not rank the rules the same way. On OLLB under \texttt{emp99} the share of feasible targets is 12.1\% for the mean (546 of 4495), 10.0\% for the median and 6.1\% for the minimum; every target is feasible up to rank 394 (mean), 90 (median) and 5 (minimum); the largest feasible rank is 901, 1207 and 2225. A bootstrap over the 540 gain pairs (1000 resamples, 95\% percentile intervals; sampling noise of the cap percentile only, not the selection bias of published fine-tunes, and not run for rank rules) gives 12.1\% [11.2, 16.6] (mean), 10.0\% [8.3, 12.1] (median) and 6.1\% [5.0, 8.2] (minimum), and 901 [824, 1201] for the mean's largest feasible rank; the mean and median intervals overlap, the minimum's lies below. The median's point estimate is below the mean's, with a shorter all-feasible prefix, but it reaches isolated targets farther down; the minimum reaches the farthest isolated target but is rarely feasible. On the 49 non-leading targets in the exactly solved set, 49 are feasible under the mean, the median and Borda, 42 under the Condorcet requirement and 36 under plurality, so these two rank rules are harder to reach than the mean; on the 500-best slice (Appendix~\ref{app:more}) the shares are 95\% (mean), 91\% (Borda), 79\% (median), 56\% (plurality), 40\% (minimum) and 30\% (Condorcet). On HELM under \texttt{emp99T} the shares are 19\% [15, 42] (mean), 33\% [21, 55] (median), 12\% [9, 21] (minimum), 21\% (Borda), 32\% (Condorcet and Copeland) and 58\% (plurality), with all-feasible prefixes of 16, 19, 6, 17, 26, 26 and 39; on HELM the point estimates of the median and Condorcet-type rules exceed the mean's, but the median and mean intervals overlap heavily, so we claim no reach ordering on HELM from this cap alone. The ordering is also a function of the cap: on OLLB the minimum's share is 2.3\%, 9.9\% and 64\% under uniform caps of $+5$, $+10$ and $+20$ points, against 0.3\%, 2.5\% and 18\% for the mean, so the minimum is the most reachable rule under these uniform caps and the least reachable under the per-task empirical caps; on HELM the median's share is at least the mean's under all five finite caps, and the minimum's is below it (Appendix Figure~\ref{fig:capgrid}). The cap, not the rule, decides feasibility far from the top: with the oracle cap \texttt{ceil} every target is feasible (Appendix Figure~\ref{fig:oracle}).

\paragraph{Cost on the same targets: typically lower than the mean's, with exceptions at the top (Q2).}
Table~\ref{tab:T3} compares each rule with the mean on the targets feasible under both with positive investment under both (we exclude targets with zero cost under either rule from the ratios; targets already first count as feasible in all reach statistics). Over all such OLLB targets the median of the per-target ratio $L_1^{\text{rule}}/L_1^{\text{mean}}$ is 0.39 for the median rule (95\% interval 0.38--0.40; 435 targets, cheaper than the mean for 99\% of them), 0.11 for the minimum (237 targets, cheaper for all), 0.27 for Borda (48), 0.52 for the Condorcet requirement (42; cheaper for 95\%) and 0.66 for plurality (36; cheaper for 78\%); the Borda, Condorcet and plurality figures cover the top 50 only. The median and minimum also touch fewer tasks (median ratio of $K$ 0.75 and 0.50), Borda the same number and the Condorcet requirement twice as many. On HELM the ratios are 0.07 (median rule, 15 targets), 0.05 (minimum, 7), 0.51 (Borda, 15), 0.21 (Condorcet and Copeland, 15) and 0.05 (plurality, 14); only 15 of the 78 targets are feasible under the mean, so the HELM statements rest on at most 15 targets. The reason is structural: the mean must close $n$ times the gap (Proposition~\ref{prop:mean}), order-statistic rules only the cheapest majority or the weakest tasks.

This ordering is a statement about typical targets, not a dominance result. At the top of OLLB it fails: the targets of mean-rank 2, 3 and 4 need 9.55, 8.68 and 8.37 points under the median rule against 4.77, 5.11 and 7.90 under the mean, and the rank-2 target needs 6.70 points under the Condorcet requirement (rank 3 and 4: 5.29, 6.25); only at rank 5 is the median far cheaper (0.70 against 23.86). The median ratio in the 2--5 stratum is $1.38$ with an interval that spans 0.03--2.0 over four targets, and the median is cheaper for one of them (Appendix Table~\ref{tab:T5}). On HELM the rank-2 target (DeepSeek v3) needs 5.58 points under Borda against 4.04 under the mean, and is already first under the minimum rule and within 0.001 point of first under the median, Condorcet and Copeland rules. A uniform capped raise of every task needs a median $L_1$ of 2.4 (median rule), 3.6 (minimum), 2.6 (Borda), 2.4 (Condorcet) and 3.0 (plurality) times the optimum on OLLB and 9.0, 38, 1.9 and 3.0 on HELM (Appendix Table~\ref{tab:T6}); for the mean the two coincide. The OLLB ordering mean $>$ Condorcet $>$ Borda $\gtrsim$ median $\gtrsim$ min in typical cost is therefore not uniform across rank bins; the ratio intervals for ranks 2--10 are wide.

\begin{table}[tp]
\caption{Median investment to become first, by stratum of the target's rank under the mean, primary caps: OLLB (\texttt{emp99}, upper block) and HELM (\texttt{emp99T}, lower block). \emph{feas.}: targets that can reach first place / targets in the stratum ($^\dagger$ OLLB Borda and Condorcet: the part of the stratum with rank $\le50$, hence 20/20 in 31--100). $K$ and $L_1$: medians over the rule's feasible targets with 95\% bootstrap intervals over targets; these are medians over different target sets for different rules and are not a ranking of the rules. $L_1/L_1^{\text{mean}}$: median [interval] of the per-target ratio over targets feasible under the rule and the mean with positive cost under both (same targets; from \texttt{paired\_*.csv}), and $n_{\text{both}}$: those targets / targets feasible under at least one of the two. ``--'': the mean itself, fewer than 3 jointly feasible targets, or no group in the analysis (HELM 2--5; see Appendix Table~\ref{tab:T5}). Generated by \texttt{figures/src/tables\_v2.py}.}
\label{tab:T3}
\centering\footnotesize
\input{../tables/T3c_rules_ollb_emp99.tex}\\[4pt]
\input{../tables/T3c_rules_helm_emp99T.tex}
\end{table}



\paragraph{Which of our predictions held.}
The plan fixed three hypotheses before the full results (H1--H3, Appendix~\ref{app:hyp}). H2 (cost rises with rank, feasibility falls as caps tighten) holds on both datasets. H1 (order-statistic and majority rules need \emph{more} tasks but possibly fewer points than the mean) holds only in part: they need fewer points on typical targets (with the exceptions above) but not more tasks, except that the Condorcet requirement touches more tasks than the mean on OLLB (ratio of $K$ about 2); on HELM the median needs fewer tasks than the mean ($K$ ratio 0.27). H3 is examined next.

\paragraph{Number of tasks (Q4).}
Appendix Figure~\ref{fig:scaling} draws random subsets of $n\in\{3,6,12,24,48\}$ HELM subjects and, in each subset, takes as the target the model ranked second or fifth \emph{under the rule being studied}, so the target model differs between rules and the curves of different rules mix the effect of the rule with that of the target. For the rank-2 target $K/n$ falls with $n$ under every rule. For the rank-5 target the mean and Borda rules keep touching about $0.4$ of the tasks (median $K$ 9 of 24 and 19 of 48 under the mean) while the median, minimum and plurality fall (median rule: $K=1$ at $n=3$ to $K=3$ at $n=48$); the mean's $L_1/n$ stays between 2.2 and 2.8 while the median rule's falls from 0.8 to 0.04. H3 (that $K/n$ of the median and majority rules does not shrink) is therefore not supported for the median on these own-rule targets. At rank 5 first place is reachable in only 55--90\% of the subsets for the minimum and 28\% (at $n=3$) for plurality, so their cost curves are conditional on feasible subsets; the bands are wide (60 draws per $n$). A common-target run (Appendix Figure~\ref{fig:scalingc}: the target is the model ranked 2nd or 5th under the \emph{mean} of each subset, one target per draw for all rules) separates the rule from the target. For the rank-5 target the median investment under the mean grows 20-fold from $n=3$ to $n=48$ (6.6 to 133 points), against 2.8-fold under the median rule (2.7 to 7.4), 1.8 under the minimum (3.6 to 6.5), 1.7 under plurality (3.5 to 5.9) and 11 under Borda (3.7 to 41), as Proposition~\ref{prop:mean} predicts for a rule that must close $n$ times the gap; $K/n$ of the median rule falls from 0.33 to 0.10, so H3 is not supported on the common target either. At rank 2 many draws already make the target first under the other rules (at $n=48$: Borda 18, median 21, minimum 32 of 60 draws), which pulls their median cost towards zero. OLLB ($n=6$) supports no scaling claim.



\paragraph{Robustness (Q5).}
Appendix Figure~\ref{fig:robust} perturbs the leaderboard and recomputes the \emph{optimal} requirement on each draw. Two perturbations are used: independent Gaussian noise on every score with the binomial standard error from the item counts and a bootstrap over tasks (draw counts and the equal-item-count assumption for OLLB are in the figure caption). What the runs measure is the variability of the optimal requirement, not whether a fixed, previously optimised raise would still win (the fixed-policy check is below). Under evaluation noise the 95\% range of $L_1$ under the mean is about $3$--$3.6$ times its median for the OLLB targets of rank 2--3 and shrinks to 0.4--0.8 times for ranks 5--100; on HELM it is $14$ times at rank 2, 2.7 at rank 3 and 0.5--0.7 at ranks 5--10. For the median rule on OLLB the ranges are 0.4--3.0 times the median, and the minimum rule is infeasible in 35--100\% of the OLLB draws. A target whose gap is within the evaluation noise therefore has no stable requirement: at ranks 2--3 some draws already make it first (lower end of the range 0). The observed requirement lies inside the 2.5--97.5 range in every feasible case shown (caption), which with ranges this wide does not bound a systematic shift. Holding the optimal raise fixed instead (\texttt{results/robust/survival\_*.csv}; for 7 OLLB and 6 HELM targets of mean-rank 2--100 and 2--40, 300 and 200 noise draws each, noise independent across cells, the policy not re-optimised; Appendix Table~\ref{tab:T7}), the target is first in a minority of draws at the smallest tested score margin or ordinal witness granularity ($10^{-3}$ points): the median over these targets of that share is 0.43 (OLLB) and 0.38 (HELM) for the mean, 0.21 and 0.05 for the median rule, 0.11 and 0.05 for the minimum and 0.19 (Borda), 0.18 (plurality) and 0.01 (Condorcet) on OLLB. A 1- and 2-point margin lifts the mean rule to 0.80 and 0.98 (OLLB) and 0.96 and 1.00 (HELM), and the OLLB median rule to 0.72 and 0.96, but not the HELM median rule (0.11, 0.28) or the minimum (0.21, 0.45 on OLLB).

\section{Discussion and limitations}
\label{sec:disc}
\paragraph{What the results say.}
Under caps taken from observed fine-tuning gains, a developer who can lift only the target's own scores, with rivals standing still, tops these leaderboards only from near the top; how near depends on the rule and the criterion. On the same targets the mean rule typically demands the most points, because it must close $n$ times the gap (Proposition~\ref{prop:mean}) whatever the spread, while the median and minimum pay only for a majority of tasks or for the weak tasks; Borda and the Condorcet requirement lie between. Reach is a different matter: on OLLB the median has a lower overall feasible share and a shorter all-feasible prefix than the mean while reaching isolated targets farther down, the minimum is cheap when feasible but needs every weak task raised and is limited by the cap, and on HELM the point estimates for the median, Condorcet-type rules and plurality exceed the mean's, with overlapping intervals for median and mean.

\paragraph{Conditions under which the conclusions fail or are not established.}
(1) \emph{The cap.} Feasibility is determined by the cap, not by the rule, far from the top (oracle rows of Appendix Table~\ref{tab:T4}), and the ordering of the rules by feasible share changes with the shape of the cap (per-task empirical against uniform). The OLLB cap uses gains of published fine-tunes over their declared base models; this is selection-biased toward successes and rests on about five tail points per task; a bootstrap over the pairs gives wide, upward-skewed intervals (OLLB mean share 11--17\%, HELM 15--42\%; Appendix Table~\ref{tab:T8}) that cover sampling noise only, not this bias, and was not run for the rank rules. The HELM cap is a transfer of the OLLB MMLU-PRO gain to 57 subjects and has not been measured on HELM. (2) \emph{Independence across tasks.} Each task is raised at no cost to others; real trade-offs would raise the cost. (3) \emph{Fixed rivals.} Others do not respond; strategic rivals need a game model such as \citet{chen2026incentives}. (4) \emph{Points, not effort.} ``Cheap'' means few score points, not little compute or data. (5) \emph{Ranks within noise.} The score margin of $10^{-3}$ points (the witness granularity for rank rules) that defines ``first'' is below the evaluation noise of the scores. The noise runs show that the optimal requirement near the top varies by multiples of its median; the fixed-policy check shows that a win by the smallest tested margin (score rules) or witness (rank rules) survives resampling in a minority of draws for every rule except the mean, where it is a coin flip (Section~\ref{sec:results}). (6) \emph{Rule coverage and size.} Rank rules were solved (under the strict model, Appendix~\ref{app:ordinal}) for the 49 best non-leading OLLB targets only (and on the 500-best slice), so their feasibility fractions there cannot be compared with those of the score rules over all targets; Copeland was solved on HELM and the 500-best slice only, and the Condorcet requirement upper-bounds it. (7) \emph{Small $n$ and $M$.} OLLB has six tasks and HELM 78 targets, of which at most 15 are feasible under the mean at the primary cap, so HELM rank bins hold few targets and their intervals are wide; scaling in $n$ comes from HELM subsets only (own-rule targets, and a common-target control with 60 draws per $n$) and is a statement about random subsets of one benchmark family; both datasets are snapshots. (8) \emph{Not done.} Kemeny, Schulze, the trimmed mean (a solver exists, no run), Copeland beyond the 500-best OLLB slice, a bootstrap of the cap for the rank rules, and survival under noise correlated across models.

\paragraph{Broader impact.}
The analysis shows leaderboard designers how exposed each rule is to targeted training. It could be read as a recipe for gaming a ranking; we think the risk is small, as first place is infeasible for most models under observed gains.

\paragraph{Reproducibility.}\label{sec:repro}
All numbers, tables and figures come from the supplementary code on pinned public data; versions, checksums, seeds, commands, run times (CPU only) and the script-to-figure/table map are in \texttt{REPRODUCE.md} and \texttt{figures/MANIFEST.md}. Only derived tables are released; no forbidden material was searched for or met (\texttt{README.md}).

\bibliographystyle{plainnat}
\bibliography{refs}

\appendix
\section{Proofs of Propositions~\ref{prop:mean}--\ref{prop:median}}
\label{app:proofs}
\paragraph{Proposition~\ref{prop:mean}.} $f(y)\ge b$ iff $\sum_t(y_t-x_t)\ge R$. Each raise $y_t-x_t$ lies in $[0,h_t]$, so the total can reach $R$ iff $\sum_th_t\ge R$, and the minimum of the (linear) total is $R$, attained by any feasible spread. Touching $k$ tasks supplies at most the sum of the $k$ largest $h_t$, which gives $K$. A raise with maximum $d$ supplies at most $\sum_t\min(d,h_t)$, so $L_\infty$ is the least $d$ with $\sum_t\min(d,h_t)\ge R$; at that $d$ the capped uniform raise supplies exactly $R$ (the sum is continuous and non-decreasing in $d$), so its total cost is $R$ as well. Without a binding cap ($d\le\min_th_t$) $\sum_t\min(d,h_t)=nd$ and $d=R/n$.
\paragraph{Proposition~\ref{prop:min}.} $\min_ty_t\ge b$ requires $y_t\ge b$ for every $t$, hence $c_t\ge b$ wherever $x_t<b$ (tasks with $x_t\ge b$ need no change). The cheapest $y_t$ on a deficient task is $b$, so $L_1=\sum_{t\in D}(b-x_t)$, $K=|D|$, and $L_\infty$ is the largest of these raises. Each is forced.
\paragraph{Proposition~\ref{prop:median}.} For odd $n$ the median is at least $b$ iff at least $m$ coordinates are. Tasks with $x_t\ge b$ already count, so $m'=(m-a)^+$ more are needed. Lifting task $t$ to $b$ costs $d_t$ and is possible iff $c_t\ge b$; costs are separable, so the cheapest set is the $m'$ smallest finite $d_t$ and $K=m'$, which cannot be smaller because each lifted task adds one to the count.

\section{Further results}
\label{app:more}
\begin{table}[h]
\caption{Closed forms (Propositions~\ref{prop:mean}--\ref{prop:median}) against the exact solvers and an independent generic LP/MILP (60-target sample), over all caps. Generated by \texttt{figures/src/tables.py} from \texttt{results/theory\_check.csv}.}
\label{tab:T2}
\centering\scriptsize
\input{../tables/T2_theory_check.tex}
\end{table}

\begin{table}[h]
\caption{Feasibility frontier under the primary cap (\texttt{emp99} on OLLB, \texttt{emp99T} on HELM) and the oracle cap (\texttt{ceil}). \emph{share}: fraction of targets for which first place is reachable (Wilson 95\% intervals); $r^*$: largest rank under the mean among feasible targets (may be an isolated case); \emph{all feas.\ up to}: largest $r$ such that every target of rank $\le r$ is feasible; \emph{first bin $<$ 1/2}: first rank bin (bins double from rank 2) in which fewer than half are feasible. $^\dagger$: the 49 non-leading targets among the 50 best under the mean only, not comparable with the score rules' fractions. Generated by \texttt{figures/src/tables.py}.}
\label{tab:T4}
\centering\scriptsize
\resizebox{\linewidth}{!}{\input{../tables/T4_frontier_ollb.tex}}\\[4pt]
\resizebox{\linewidth}{!}{\input{../tables/T4_frontier_helm.tex}}
\end{table}
\begin{figure}[h]
\centering
\includegraphics[width=\linewidth]{F_robust.pdf}
\caption{Stability of the optimal requirement ($L_1$) under evaluation noise (top row) and a bootstrap over tasks (bottom row), for target ranks 2, 3, 5, 10, 20, 50, 100 (OLLB, \texttt{emp99}) and 2, 3, 5, 10, 20, 40 (HELM, \texttt{emp99T}), per rule. Dot: median over feasible draws; bar: 2.5--97.5 percentile of those draws (a range of draws, not a confidence interval of a median); open diamond: value on the observed data (absent where the target is outside the solved ordinal set or infeasible; rank-rule $L_1$ is the witness cost, within $n\varepsilon$ of the strict infimum); lower strip: percentage of draws in which first place is unreachable. Noise: independent Gaussian draws with the binomial standard error of each score (300 draws for OLLB closed forms, 60 for OLLB rank rules, 200 for HELM; the task bootstrap uses the same counts; OLLB item counts come from one results file and are assumed equal for all models). Each draw is re-optimised, so the figure shows variability of the \emph{optimal requirement}, not whether a fixed raise remains a win. At ranks 2--3 the lower bar end is 0 because some draws already make the target first. Counted for the evaluation-noise row (feasible target-by-rule cases of the five rules shown), the observed cost lies inside the 2.5--97.5 range in every case (OLLB closed forms 18/18, OLLB rank rules 12/12, HELM 21/21); with ranges this wide this does not bound a shift. Generated by \texttt{figures/src/fig\_robust.py}.}
\label{fig:robust}
\end{figure}
\begin{figure}[h]
\centering
\includegraphics[width=\linewidth]{F_ollb_oracle.pdf}\\[2pt]
\includegraphics[width=\linewidth]{F_helm_oracle.pdf}
\caption{As Figure~\ref{fig:main} with the oracle cap \texttt{ceil} (scores up to 100): a lower bound on any required investment; every target is feasible, so the differences between rules are not caused by the cap.}
\label{fig:oracle}
\end{figure}
\begin{figure}[h]
\centering
\includegraphics[width=\linewidth]{F_ollb_capgrid.pdf}
\caption{OLLB: feasibility (left) and median $L_1/n$ among feasible targets (right) for three rank strata under the mean (mean-ranks 2--5, 6--20, 21--100; $N=4$ in the first, so its intervals are wide), under each cap \emph{scenario}. The horizontal axis orders scenarios and is not a monotone sweep: $+h$ is a uniform cap, \texttt{emp95} and \texttt{emp99} are per-task percentile caps, ``front'' is the best rival score on each task, 100 is the ceiling; these cap vectors are not nested. Two things move the right panels: the feasible cohort changes with the cap, so a median over feasible targets changes even if no target's own optimum does; and a target's own optimum, which under the mean (and for the minimum, when feasible) is the same under any cap under which it is feasible (Proposition~\ref{prop:mean}), while for the median and rank rules no such invariance is established and relaxing a cap can allow a cheaper allocation. Where the mean line is flat the cohort did not change; where it moves (6--20 at $+5$) it did. Bands: Wilson 95\% (left), bootstrap over targets (right), absent with fewer than 3 feasible targets. Borda and Condorcet (50 best targets only) are unconnected markers and not comparable with the score rules' fractions in the 21--100 stratum. Generated by \texttt{figures/src/fig\_cap\_grid.py}.}
\label{fig:capgrid}
\end{figure}
\begin{figure}[h]
\centering
\includegraphics[width=\linewidth]{F_helm_capgrid.pdf}
\caption{HELM: as Figure~\ref{fig:capgrid} (strata 2--5, 6--20, 21--78); $T$ denotes the transfer of the OLLB MMLU-PRO gains to all 57 subjects, so the scenarios mix a transfer assumption with the cap shape.}
\label{fig:capgridh}
\end{figure}

\begin{table}[h]
\caption{The four non-leading targets of mean-rank 2--5: minimum $L_1$ in points ($K$ in parentheses) under each rule; OLLB (\texttt{emp99}, upper) and HELM (\texttt{emp99T}, lower). Bold: cheapest rule(s) for that target; ``0.00 (1)'' is a raise of the 0.001 margin, ``0.00 (0)'' means the target is already first under that rule. Last row: number of the four targets for which the rule is strictly cheaper than the mean. With four targets per table these are exceptions to a general ordering, not rates. Generated by \texttt{figures/src/tables\_v2.py}.}
\label{tab:T5}
\centering\footnotesize
\input{../tables/T5_exceptions_ollb.tex}\\[4pt]
\input{../tables/T5_exceptions_helm.tex}
\end{table}
\begin{figure}[h]
\centering
\includegraphics[width=\linewidth]{F_helm_scaling.pdf}
\caption{HELM MMLU: random subsets of $n$ of the 57 subjects (60 subsets per $n$), cap \texttt{emp99T}. In every subset and for every rule separately, the target is the model ranked 2 (left) or 5 (right) \emph{under that rule on that subset}, so the target differs between rules and subsets and the curves mix the effect of the rule with that of the target; it is not the mean-ranked model. Top row: percentage of subsets in which first place is reachable within the cap (counts in \texttt{F\_helm\_scaling.json}). Middle and bottom rows: $K/n$ and $L_1/n$, medians over the \emph{feasible} subsets only; bands are the 2.5--97.5 percentile range of those subsets (a range of draws, not a confidence interval of the median). Where the top row is below 100\% the cost curve is conditional on a selected set of subsets. The Condorcet requirement is not in this run. Generated by \texttt{figures/src/fig\_scaling.py}.}
\label{fig:scaling}
\end{figure}

\begin{figure}[h]
\centering
\includegraphics[width=\linewidth]{F_helm_scaling_common.pdf}
\caption{HELM scaling control: as Figure~\ref{fig:scaling}, but the target is the same model under every rule within a draw (rank 2 or 5 under the \emph{mean} of the subset; 60 draws per $n$, cap \texttt{emp99T}). Rows: share of draws in which first place is reachable; share in which the target is already first (cost 0, included in the cost rows); median $K/n$ and $L_1/n$ over feasible draws with the 2.5--97.5 percentile range of those draws (a spread of draws, not an interval for the median); the $L_1/n$ axis is symlog. Only this version supports a statement about how a rule reacts to $n$. Generated by \texttt{figures/src/fig\_scaling\_common.py}.}
\label{fig:scalingc}
\end{figure}

\begin{table}[h]
\caption{Capped uniform raise of every task by the same amount against the optimal raise, primary cap (OLLB \texttt{emp99}, upper; HELM \texttt{emp99T}, lower): ratio $L_1^{\text{unif}}/L_1^{\text{opt}}$, median over targets feasible under both with positive cost [95\% bootstrap interval over targets]; \emph{ratio $<1$}: targets with a ratio below 1 (rank rules only, at most $n\varepsilon/L_1^{\text{opt}}$ because the optimum is a witness cost); $K$: medians. $^\dagger$: OLLB rank rules cover the 50 best targets only. The sets of feasible targets differ by rule. Generated by \texttt{figures/src/tables\_v3.py}.}
\label{tab:T6}
\centering\scriptsize
\input{../tables/T6_uniform_ollb.tex}\\[4pt]
\input{../tables/T6_uniform_helm.tex}
\end{table}

\begin{table}[h]
\caption{Survival of the optimised raise under evaluation noise, the raise fixed (not re-optimised): share of noise draws in which the target is still strictly first; median over targets, min--max over targets, and the median with no raise. $\varepsilon$ is the score margin (0.001, the smallest used, 1 or 2 points) for the mean, median and minimum, and the witness granularity (0.001) for the rank rules. OLLB: 7 targets of mean-rank 2--100, 300 draws each; HELM: 6 targets of rank 2--40, 200 draws each; infeasible and already-first targets are excluded, so the target count falls under some rules. Medians over 2--7 targets are descriptive. Generated by \texttt{figures/src/tables\_v3.py}.}
\label{tab:T7}
\centering\scriptsize
\input{../tables/T7_survival.tex}
\end{table}

\begin{table}[h]
\caption{Cap bootstrap (1000 resamples of the 540 fine-tune/base pairs, each pair drawn as a whole row with all its tasks so the dependence across tasks is kept; HELM: the transferred OLLB MMLU-PRO gains): feasible share, largest feasible rank under the mean and all-feasible prefix, observed value with the 2.5--97.5 percentile interval. The intervals cover sampling noise of the cap only, not the selection bias of published fine-tunes or the HELM transfer; rank rules were not bootstrapped. Generated by \texttt{figures/src/tables\_v3.py}.}
\label{tab:T8}
\centering\scriptsize
\input{../tables/T8_capboot.tex}
\end{table}

\subsection*{Variants and sensitivity checks}
All numbers below are in \texttt{results/summary/summary\_variants.csv} (script \texttt{code/summarize\_variants.py}). Ratios are over targets feasible both in the variant and in the core run.
\paragraph{Margin.} With margins of 1 and 2 points instead of $10^{-3}$ the feasible share on OLLB falls from 12.1\% to 11.1\% and 10.4\% (mean), from 10.0\% to 9.0\% and 7.6\% (median) and from 6.1\% to 4.3\% and 2.6\% (minimum); the median ratio of the cost to the minimal-margin cost is 1.08 and 1.17 (mean), 1.15 and 1.34 (median), 1.25 and 1.53 (minimum). On HELM the shares fall from 19.2\% to 15.4\% and 11.5\% (mean), from 33.3\% to 28.2\% and 20.5\% (median), and from 11.5\% to 11.5\% and 9.0\% (minimum), with ratios of 1.3 and 1.7 (mean), 1.5 and 2.3 (median). Margin variants exist for the score rules only: the rank rules use strict comparisons without a margin (Appendix~\ref{app:ordinal}). The ordering of the feasible shares is unchanged on OLLB.
\paragraph{Slice of the 500 best models.} To extend the rank rules beyond the 50 best targets we solved them on a leaderboard that keeps only the 500 best models by the mean (499 non-leading targets; time-outs are counted as censored, not as infeasible, and none remain under \texttt{emp99} after a 900\,s re-solve). Under \texttt{emp99} the feasible shares are 94.8\% (mean), 90.6\% (Borda), 79.4\% (median), 74.9\% (Copeland), 56.3\% (plurality), 40.3\% (minimum) and 30.5\% (Condorcet requirement, which is sufficient for Copeland and hence below it); with the oracle cap every target is feasible under every rule (Borda: 99.8\% with one censored instance). These shares refer to a different set of rivals than the full leaderboard, and the rules are again not compared on a common footing beyond what the paired columns of \texttt{summary\_variants.csv} show.
\paragraph{Merged models, scale and weights.} Dropping the merged models (3788 targets) leaves the OLLB shares nearly unchanged (mean 11.5\%, median 9.4\%, minimum 7.0\%). With raw task accuracies instead of the leaderboard's normalised scores the ordering of the feasible shares under a +10 point cap changes (normalised: minimum 9.9\%, median 4.5\%, mean 2.5\%; raw: mean 6.6\%, minimum 5.0\%, median 0.5\%), so cap comparisons in points depend on the score scale; we use the leaderboard's own normalisation. Random task weights drawn from a Dirichlet distribution (concentration 1 or 5; the stratified targets of ranks 2, 3, 5, 10, 20, 50, 100 on OLLB and 2, 3, 5, 10, 20, 40 on HELM, 500 draws each) give a median weighted cost of 0.54 and 0.70 (OLLB) and 0.43 and 0.68 (HELM) times the equal-weight cost of the same target, so weights that concentrate on tasks where the target is strong make first place cheaper.
\paragraph{Further score rules.} Geometric, harmonic and squared power means and plurality were solved on the same targets; their results are in \texttt{results/core\_*.csv} and the frontier tables \texttt{results/summary/frontier\_*.csv}, and we make no separate claim about them in the text. A trimmed mean has a solver in the code but was not run.

\section{Exactness of the rank-rule solver}
\label{app:ordinal}
The rank rules (Borda, Condorcet requirement, plurality, Copeland) depend only on the order of scores on each task. We therefore use one model for every cost and for the winner test. A pairwise comparison of two scores is a win (strictly higher), a loss, or a tie worth one half; Borda adds, per task, the models strictly below plus half the ties; the Condorcet requirement asks for strictly more than $n/2$ pairwise points against every rival; Copeland scores each rival beaten as 1 and tied as one half; plurality gives each task one point split equally among all $m$ models tied for its highest score (so a tie is worth $1/m$, not one half). Under Borda, plurality and Copeland the target is first if its total is strictly the largest; under the Condorcet requirement it is first if it has more than $n/2$ pairwise points against each rival. There is no margin. A strict win is determined by order, so the strict-winning infimum of $L_1$ need not be attained; it is attained when a tie suffices (for Borda, a target with scores $(0,1)$ against rivals $(1,0)$ and $(1,0.5)$, cap $(1,1)$: any task-0 score below 1 leaves it behind the second rival, a score of 1 ties both rivals on that task and wins, so the minimum $L_1=1$ is attained at a tie). The solver restricts each raised task to a rival's score (a tie) or to $\min(s+\varepsilon,c_t)$ with $c_t>s$ (strictly above a rival score $s$, at most $\varepsilon$ above it), and returns the cost of that explicit witness. Two consequences follow. Feasibility and $K$ do not depend on $\varepsilon$ and are those of continuous raises; $L_\infty$ is the infimum itself (bisection, tolerance $10^{-6}$). The infimum of $L_1$ lies in $[L_1-n\varepsilon,L_1]$ ($n$ tasks): by monotonicity, any winning continuous raise can be replaced by an allowed witness that costs at most $\varepsilon$ more per raised task, so the witness optimum is at most the infimum plus $n\varepsilon$, and it is never below the infimum. The bound uses $n\varepsilon$ and not $K\varepsilon$, because the $K$ reported beside $L_1$ is minimised separately and need not be the number of tasks raised in the witness. An earlier version used finite levels at $s+\varepsilon$ credited with passing every lower rival; a review of the code showed that this model could disagree with $L_\infty$ and with the pairwise-margin definition on closely spaced scores (a level generated from one rival could pass another by less than $\varepsilon$ and be credited in full), and the paper's earlier rank-rule costs came from it. The strict model replaced it, with plurality solved on an exact credit lattice and an exact rational check of every witness. The unit tests compare the solver with brute force on fractional, tied, sub-$\varepsilon$ and cap-boundary instances and on the reviewers' counterexamples. We regenerated every rank-rule result with it (\texttt{results/audit\_ordinal\_fix.csv}): no feasibility verdict changed in any file, only Borda costs changed in the primary files (4 of 100 rows of the OLLB top-50 file, largest change 0.10 points; none on HELM), and more in the sensitivity runs (Borda on the 500-best slice 116 of 1000 rows, up to 1.35 points; HELM cap sweep 9 of 395 rows, up to 3.5 points; OLLB task bootstrap 20 of 420 rows, up to 1.6 points). All numbers in this paper come from the regenerated files. For score rules a target already ahead of every rival costs zero, whatever the size of its lead; this matters for the margin variants only for the leaders of the median and minimum rules among the mean-nonleaders. $L_\infty$ is not reported for rank rules in the main text, and the rank-rule noise and bootstrap rows are over the same strict model.

\section{Hypotheses fixed before the full results}
\label{app:hyp}
The plan fixed three hypotheses from the closed forms before the full result set was available. H1: with loose caps, score-sum rules need the same $L_1$ regardless of $K$, while order-statistic and majority rules need more tasks but possibly fewer points than the mean. H2: required investment grows with the target's initial rank; feasibility falls as caps tighten; rule ordering can change with the cap. H3: for fixed per-task caps the share $K/n$ of tasks that majority and median rules need does not shrink with $n$, while the mean rule's $L_1/n$ depends on the density of gaps. Anything beyond these is exploratory.

\clearpage
\section*{NeurIPS Paper Checklist}




\begin{enumerate}

\item {\bf Claims}
    \item[] Question: Do the main claims made in the abstract and introduction accurately reflect the paper's contributions and scope?
    \item[] Answer: \answerYes{}
    \item[] Justification: Abstract and introduction state the two dataset-specific findings, the cap assumption and their scope; each is backed by Table~\ref{tab:T3}, Appendix Table~\ref{tab:T4} and Figure~\ref{fig:main}, and Section~\ref{sec:disc} lists where they fail.
    \item[] Guidelines:
    \begin{itemize}
        \item The answer \answerNA{} means that the abstract and introduction do not include the claims made in the paper.
        \item The abstract and/or introduction should clearly state the claims made, including the contributions made in the paper and important assumptions and limitations. A \answerNo{} or \answerNA{} answer to this question will not be perceived well by the reviewers. 
        \item The claims made should match theoretical and experimental results, and reflect how much the results can be expected to generalize to other settings. 
        \item It is fine to include aspirational goals as motivation as long as it is clear that these goals are not attained by the paper. 
    \end{itemize}

\item {\bf Limitations}
    \item[] Question: Does the paper discuss the limitations of the work performed by the authors?
    \item[] Answer: \answerYes{}
    \item[] Justification: Section~\ref{sec:disc} (Discussion and limitations) lists cap selection bias and transfer, task independence, fixed rivals, noise, rule coverage, snapshots and unfinished items.
    \item[] Guidelines:
    \begin{itemize}
        \item The answer \answerNA{} means that the paper has no limitation while the answer \answerNo{} means that the paper has limitations, but those are not discussed in the paper. 
        \item The authors are encouraged to create a separate ``Limitations'' section in their paper.
        \item The paper should point out any strong assumptions and how robust the results are to violations of these assumptions (e.g., independence assumptions, noiseless settings, model well-specification, asymptotic approximations only holding locally). The authors should reflect on how these assumptions might be violated in practice and what the implications would be.
        \item The authors should reflect on the scope of the claims made, e.g., if the approach was only tested on a few datasets or with a few runs. In general, empirical results often depend on implicit assumptions, which should be articulated.
        \item The authors should reflect on the factors that influence the performance of the approach. For example, a facial recognition algorithm may perform poorly when image resolution is low or images are taken in low lighting. Or a speech-to-text system might not be used reliably to provide closed captions for online lectures because it fails to handle technical jargon.
        \item The authors should discuss the computational efficiency of the proposed algorithms and how they scale with dataset size.
        \item If applicable, the authors should discuss possible limitations of their approach to address problems of privacy and fairness.
        \item While the authors might fear that complete honesty about limitations might be used by reviewers as grounds for rejection, a worse outcome might be that reviewers discover limitations that aren't acknowledged in the paper. The authors should use their best judgment and recognize that individual actions in favor of transparency play an important role in developing norms that preserve the integrity of the community. Reviewers will be specifically instructed to not penalize honesty concerning limitations.
    \end{itemize}

\item {\bf Theory assumptions and proofs}
    \item[] Question: For each theoretical result, does the paper provide the full set of assumptions and a complete (and correct) proof?
    \item[] Answer: \answerYes{}
    \item[] Justification: Propositions~\ref{prop:mean}--\ref{prop:median} state their assumptions in Section~\ref{sec:solvers} and are proved in Appendix~\ref{app:proofs}; rank-rule results are computational (witness costs within $n\varepsilon$ of the infimum of a strict model, Appendix~\ref{app:ordinal}) and no theorem is claimed for them.
    \item[] Guidelines:
    \begin{itemize}
        \item The answer \answerNA{} means that the paper does not include theoretical results. 
        \item All the theorems, formulas, and proofs in the paper should be numbered and cross-referenced.
        \item All assumptions should be clearly stated or referenced in the statement of any theorems.
        \item The proofs can either appear in the main paper or the supplemental material, but if they appear in the supplemental material, the authors are encouraged to provide a short proof sketch to provide intuition. 
        \item Inversely, any informal proof provided in the core of the paper should be complemented by formal proofs provided in appendix or supplemental material.
        \item Theorems and Lemmas that the proof relies upon should be properly referenced. 
    \end{itemize}

    \item {\bf Experimental result reproducibility}
    \item[] Question: Does the paper fully disclose all the information needed to reproduce the main experimental results of the paper to the extent that it affects the main claims and/or conclusions of the paper (regardless of whether the code and data are provided or not)?
    \item[] Answer: \answerYes{}
    \item[] Justification: Section~\ref{sec:setup}--\ref{sec:exp} define rules, caps, margin and filters; \texttt{REPRODUCE.md} in the supplementary code gives data versions, seeds and commands.
    \item[] Guidelines:
    \begin{itemize}
        \item The answer \answerNA{} means that the paper does not include experiments.
        \item If the paper includes experiments, a \answerNo{} answer to this question will not be perceived well by the reviewers: Making the paper reproducible is important, regardless of whether the code and data are provided or not.
        \item If the contribution is a dataset and\slash or model, the authors should describe the steps taken to make their results reproducible or verifiable. 
        \item Depending on the contribution, reproducibility can be accomplished in various ways. For example, if the contribution is a novel architecture, describing the architecture fully might suffice, or if the contribution is a specific model and empirical evaluation, it may be necessary to either make it possible for others to replicate the model with the same dataset, or provide access to the model. In general. releasing code and data is often one good way to accomplish this, but reproducibility can also be provided via detailed instructions for how to replicate the results, access to a hosted model (e.g., in the case of a large language model), releasing of a model checkpoint, or other means that are appropriate to the research performed.
        \item While NeurIPS does not require releasing code, the conference does require all submissions to provide some reasonable avenue for reproducibility, which may depend on the nature of the contribution. For example
        \begin{enumerate}
            \item If the contribution is primarily a new algorithm, the paper should make it clear how to reproduce that algorithm.
            \item If the contribution is primarily a new model architecture, the paper should describe the architecture clearly and fully.
            \item If the contribution is a new model (e.g., a large language model), then there should either be a way to access this model for reproducing the results or a way to reproduce the model (e.g., with an open-source dataset or instructions for how to construct the dataset).
            \item We recognize that reproducibility may be tricky in some cases, in which case authors are welcome to describe the particular way they provide for reproducibility. In the case of closed-source models, it may be that access to the model is limited in some way (e.g., to registered users), but it should be possible for other researchers to have some path to reproducing or verifying the results.
        \end{enumerate}
    \end{itemize}


\item {\bf Open access to data and code}
    \item[] Question: Does the paper provide open access to the data and code, with sufficient instructions to faithfully reproduce the main experimental results, as described in supplemental material?
    \item[] Answer: \answerYes{}
    \item[] Justification: The data are public (Open LLM Leaderboard and HELM, versions and checksums in \texttt{REPRODUCE.md}); the supplementary code reproduces every table and figure with the commands listed there.
    \item[] Guidelines:
    \begin{itemize}
        \item The answer \answerNA{} means that paper does not include experiments requiring code.
        \item Please see the NeurIPS code and data submission guidelines (\url{https://neurips.cc/public/guides/CodeSubmissionPolicy}) for more details.
        \item While we encourage the release of code and data, we understand that this might not be possible, so \answerNo{} is an acceptable answer. Papers cannot be rejected simply for not including code, unless this is central to the contribution (e.g., for a new open-source benchmark).
        \item The instructions should contain the exact command and environment needed to run to reproduce the results. See the NeurIPS code and data submission guidelines (\url{https://neurips.cc/public/guides/CodeSubmissionPolicy}) for more details.
        \item The authors should provide instructions on data access and preparation, including how to access the raw data, preprocessed data, intermediate data, and generated data, etc.
        \item The authors should provide scripts to reproduce all experimental results for the new proposed method and baselines. If only a subset of experiments are reproducible, they should state which ones are omitted from the script and why.
        \item At submission time, to preserve anonymity, the authors should release anonymized versions (if applicable).
        \item Providing as much information as possible in supplemental material (appended to the paper) is recommended, but including URLs to data and code is permitted.
    \end{itemize}


\item {\bf Experimental setting/details}
    \item[] Question: Does the paper specify all the training and test details (e.g., data splits, hyperparameters, how they were chosen, type of optimizer) necessary to understand the results?
    \item[] Answer: \answerYes{}
    \item[] Justification: Section~\ref{sec:exp} and Appendix state data filters, dedupe, caps, bootstrap and Monte Carlo settings; \texttt{REPRODUCE.md} lists the rest.
    \item[] Guidelines:
    \begin{itemize}
        \item The answer \answerNA{} means that the paper does not include experiments.
        \item The experimental setting should be presented in the core of the paper to a level of detail that is necessary to appreciate the results and make sense of them.
        \item The full details can be provided either with the code, in appendix, or as supplemental material.
    \end{itemize}

\item {\bf Experiment statistical significance}
    \item[] Question: Does the paper report error bars suitably and correctly defined or other appropriate information about the statistical significance of the experiments?
    \item[] Answer: \answerYes{}
    \item[] Justification: Bootstrap intervals over targets (Figure~\ref{fig:main}, Table~\ref{tab:T3}), Wilson intervals for feasible shares (Appendix Table~\ref{tab:T4}), a bootstrap over the gain pairs for the cap percentile (feasible shares and frontier ranks, Section~\ref{sec:results}), Monte Carlo ranges for evaluation noise and task resampling on both datasets (Figure~\ref{fig:robust}), and the share of draws in which a fixed raise still wins, stated as ranges or medians over draws. Not quantified: the cap bootstrap for rank rules and noise correlated across models (Section~\ref{sec:disc}).
    \item[] Guidelines:
    \begin{itemize}
        \item The answer \answerNA{} means that the paper does not include experiments.
        \item The authors should answer \answerYes{} if the results are accompanied by error bars, confidence intervals, or statistical significance tests, at least for the experiments that support the main claims of the paper.
        \item The factors of variability that the error bars are capturing should be clearly stated (for example, train/test split, initialization, random drawing of some parameter, or overall run with given experimental conditions).
        \item The method for calculating the error bars should be explained (closed form formula, call to a library function, bootstrap, etc.)
        \item The assumptions made should be given (e.g., Normally distributed errors).
        \item It should be clear whether the error bar is the standard deviation or the standard error of the mean.
        \item It is OK to report 1-sigma error bars, but one should state it. The authors should preferably report a 2-sigma error bar than state that they have a 96\% CI, if the hypothesis of Normality of errors is not verified.
        \item For asymmetric distributions, the authors should be careful not to show in tables or figures symmetric error bars that would yield results that are out of range (e.g., negative error rates).
        \item If error bars are reported in tables or plots, the authors should explain in the text how they were calculated and reference the corresponding figures or tables in the text.
    \end{itemize}

\item {\bf Experiments compute resources}
    \item[] Question: For each experiment, does the paper provide sufficient information on the computer resources (type of compute workers, memory, time of execution) needed to reproduce the experiments?
    \item[] Answer: \answerYes{}
    \item[] Justification: The Reproducibility paragraph of Section~\ref{sec:disc} states CPU only; \texttt{REPRODUCE.md} gives the CPU type, memory and run times per step.
    \item[] Guidelines:
    \begin{itemize}
        \item The answer \answerNA{} means that the paper does not include experiments.
        \item The paper should indicate the type of compute workers CPU or GPU, internal cluster, or cloud provider, including relevant memory and storage.
        \item The paper should provide the amount of compute required for each of the individual experimental runs as well as estimate the total compute. 
        \item The paper should disclose whether the full research project required more compute than the experiments reported in the paper (e.g., preliminary or failed experiments that didn't make it into the paper). 
    \end{itemize}
    
\item {\bf Code of ethics}
    \item[] Question: Does the research conducted in the paper conform, in every respect, with the NeurIPS Code of Ethics \url{https://neurips.cc/public/EthicsGuidelines}?
    \item[] Answer: \answerYes{}
    \item[] Justification: The work uses only public score tables, trains no models and involves no human subjects; we have read the NeurIPS Code of Ethics and see no conflict.
    \item[] Guidelines:
    \begin{itemize}
        \item The answer \answerNA{} means that the authors have not reviewed the NeurIPS Code of Ethics.
        \item If the authors answer \answerNo, they should explain the special circumstances that require a deviation from the Code of Ethics.
        \item The authors should make sure to preserve anonymity (e.g., if there is a special consideration due to laws or regulations in their jurisdiction).
    \end{itemize}


\item {\bf Broader impacts}
    \item[] Question: Does the paper discuss both potential positive societal impacts and negative societal impacts of the work performed?
    \item[] Answer: \answerYes{}
    \item[] Justification: Section~\ref{sec:disc} (Broader impact) discusses benefits for leaderboard designers and the risk that the analysis is read as a manipulation recipe.
    \item[] Guidelines:
    \begin{itemize}
        \item The answer \answerNA{} means that there is no societal impact of the work performed.
        \item If the authors answer \answerNA{} or \answerNo, they should explain why their work has no societal impact or why the paper does not address societal impact.
        \item Examples of negative societal impacts include potential malicious or unintended uses (e.g., disinformation, generating fake profiles, surveillance), fairness considerations (e.g., deployment of technologies that could make decisions that unfairly impact specific groups), privacy considerations, and security considerations.
        \item The conference expects that many papers will be foundational research and not tied to particular applications, let alone deployments. However, if there is a direct path to any negative applications, the authors should point it out. For example, it is legitimate to point out that an improvement in the quality of generative models could be used to generate Deepfakes for disinformation. On the other hand, it is not needed to point out that a generic algorithm for optimizing neural networks could enable people to train models that generate Deepfakes faster.
        \item The authors should consider possible harms that could arise when the technology is being used as intended and functioning correctly, harms that could arise when the technology is being used as intended but gives incorrect results, and harms following from (intentional or unintentional) misuse of the technology.
        \item If there are negative societal impacts, the authors could also discuss possible mitigation strategies (e.g., gated release of models, providing defenses in addition to attacks, mechanisms for monitoring misuse, mechanisms to monitor how a system learns from feedback over time, improving the efficiency and accessibility of ML).
    \end{itemize}
    
\item {\bf Safeguards}
    \item[] Question: Does the paper describe safeguards that have been put in place for responsible release of data or models that have a high risk for misuse (e.g., pre-trained language models, image generators, or scraped datasets)?
    \item[] Answer: \answerNA{}
    \item[] Justification: No models or data with high misuse risk are released; only derived score tables and code.
    \item[] Guidelines:
    \begin{itemize}
        \item The answer \answerNA{} means that the paper poses no such risks.
        \item Released models that have a high risk for misuse or dual-use should be released with necessary safeguards to allow for controlled use of the model, for example by requiring that users adhere to usage guidelines or restrictions to access the model or implementing safety filters. 
        \item Datasets that have been scraped from the Internet could pose safety risks. The authors should describe how they avoided releasing unsafe images.
        \item We recognize that providing effective safeguards is challenging, and many papers do not require this, but we encourage authors to take this into account and make a best faith effort.
    \end{itemize}

\item {\bf Licenses for existing assets}
    \item[] Question: Are the creators or original owners of assets (e.g., code, data, models), used in the paper, properly credited and are the license and terms of use explicitly mentioned and properly respected?
    \item[] Answer: \answerYes{}
    \item[] Justification: Sources and licence status are listed in \texttt{REPRODUCE.md}: neither leaderboard declares a licence for its score data, HELM code is Apache-2.0; we say so and release only derived tables.
    \item[] Guidelines:
    \begin{itemize}
        \item The answer \answerNA{} means that the paper does not use existing assets.
        \item The authors should cite the original paper that produced the code package or dataset.
        \item The authors should state which version of the asset is used and, if possible, include a URL.
        \item The name of the license (e.g., CC-BY 4.0) should be included for each asset.
        \item For scraped data from a particular source (e.g., website), the copyright and terms of service of that source should be provided.
        \item If assets are released, the license, copyright information, and terms of use in the package should be provided. For popular datasets, \url{paperswithcode.com/datasets} has curated licenses for some datasets. Their licensing guide can help determine the license of a dataset.
        \item For existing datasets that are re-packaged, both the original license and the license of the derived asset (if it has changed) should be provided.
        \item If this information is not available online, the authors are encouraged to reach out to the asset's creators.
    \end{itemize}

\item {\bf New assets}
    \item[] Question: Are new assets introduced in the paper well documented and is the documentation provided alongside the assets?
    \item[] Answer: \answerYes{}
    \item[] Justification: The released code and derived score tables are documented in \texttt{REPRODUCE.md}; no other new assets are released.
    \item[] Guidelines:
    \begin{itemize}
        \item The answer \answerNA{} means that the paper does not release new assets.
        \item Researchers should communicate the details of the dataset\slash code\slash model as part of their submissions via structured templates. This includes details about training, license, limitations, etc. 
        \item The paper should discuss whether and how consent was obtained from people whose asset is used.
        \item At submission time, remember to anonymize your assets (if applicable). You can either create an anonymized URL or include an anonymized zip file.
    \end{itemize}

\item {\bf Crowdsourcing and research with human subjects}
    \item[] Question: For crowdsourcing experiments and research with human subjects, does the paper include the full text of instructions given to participants and screenshots, if applicable, as well as details about compensation (if any)? 
    \item[] Answer: \answerNA{}
    \item[] Justification: The paper involves no crowdsourcing or human subjects.
    \item[] Guidelines:
    \begin{itemize}
        \item The answer \answerNA{} means that the paper does not involve crowdsourcing nor research with human subjects.
        \item Including this information in the supplemental material is fine, but if the main contribution of the paper involves human subjects, then as much detail as possible should be included in the main paper. 
        \item According to the NeurIPS Code of Ethics, workers involved in data collection, curation, or other labor should be paid at least the minimum wage in the country of the data collector. 
    \end{itemize}

\item {\bf Institutional review board (IRB) approvals or equivalent for research with human subjects}
    \item[] Question: Does the paper describe potential risks incurred by study participants, whether such risks were disclosed to the subjects, and whether Institutional Review Board (IRB) approvals (or an equivalent approval/review based on the requirements of your country or institution) were obtained?
    \item[] Answer: \answerNA{}
    \item[] Justification: The paper involves no human subjects, so no IRB approval applies.
    \item[] Guidelines:
    \begin{itemize}
        \item The answer \answerNA{} means that the paper does not involve crowdsourcing nor research with human subjects.
        \item Depending on the country in which research is conducted, IRB approval (or equivalent) may be required for any human subjects research. If you obtained IRB approval, you should clearly state this in the paper. 
        \item We recognize that the procedures for this may vary significantly between institutions and locations, and we expect authors to adhere to the NeurIPS Code of Ethics and the guidelines for their institution. 
        \item For initial submissions, do not include any information that would break anonymity (if applicable), such as the institution conducting the review.
    \end{itemize}

\item {\bf Declaration of LLM usage}
    \item[] Question: Does the paper describe the usage of LLMs if it is an important, original, or non-standard component of the core methods in this research? Note that if the LLM is used only for writing, editing, or formatting purposes and does \emph{not} impact the core methodology, scientific rigor, or originality of the research, declaration is not required.
    %this research? 
    \item[] Answer: \answerNA{}
    \item[] Justification: No LLM is part of the method; LLM assistance was used only for writing and coding assistance.
    \item[] Guidelines:
    \begin{itemize}
        \item The answer \answerNA{} means that the core method development in this research does not involve LLMs as any important, original, or non-standard components.
        \item Please refer to our LLM policy in the NeurIPS handbook for what should or should not be described.
    \end{itemize}

\end{enumerate}

\end{document}
